% Appendix: project start to 7 Oct 2026, and every reproduction of published work. Draft 27 Nov 2026 (writer id: repro).
% Every number and dated fact is traced in app_repro_NOTES.md. Old-era numbers are history only (RECORD_BRIEF rule 4).
\chapter{Reproducing published work and the first pipeline}\label{app:repro}

This appendix records the project from its start to the audit and restart of 7 and 8 Oct 2026, and every reproduction of
published work we ran, including two that ran later (13 and 14 Oct). Chapter~\ref{ch:related} reviews the published numbers and why
they do not compare, Tables~\ref{tab:T13a} and~\ref{tab:T13b} re-score our final models in the published conventions, and
Chapter~\ref{ch:legacy} continues the history from 10 Oct. Nothing here decides anything in the thesis.

\paragraph{Conditions of every old-era number.} All numbers from before 7 Oct are history and reproduction records only. On BenchmarkEL, unless
stated otherwise, they come from the first pipeline: the original BenchmarkEL \texttt{20211104} masks, a U-Net++ trained from scratch, class-weighted
cross-entropy (CE), labelled batch 4, a cosine learning rate that ends at zero, and the final step of 20,000 steps. The 50 legacy test
cells were scored in every run and were the primary readout of every comparison. The official split shares modules between train and
test (24 of the 50 test cells came from modules with train cells, as counted then; the later audit found 8 test cells that are the
same physical cell as a cell of the training pool, Section~\ref{sec:data-split}), and the class weights of the first code were
computed over all 582 cells, test included. Phase~0 numbers survive only in notes, because the code was lost, and were never
re-verified against files. Mean Teacher results of the period used a loss later shown to be invalid (Section~\ref{sec:repro-audit}).
No old-era number is a valid result, and none may be set beside Chapters~\ref{ch:protocol} to~\ref{ch:t3b}.

Unless stated otherwise, ``standard'' IoU is pooled: one dataset-level confusion matrix over all pixels of the scored set, with
background inside the matrix, so that false positives on background count. ``Defect mIoU'' is the mean of crack and inactive IoU on
BenchmarkEL and the mean over the four defect classes on UCF-EL, both in that convention. Published conventions are named where they
appear.

\section{The project start and phase 0 (to the end of September 2026)}\label{sec:repro-phase0}

\paragraph{Framing.} The project began as a master's project with the industry partner on semi-supervised (SSL) segmentation of
defects in electroluminescence (EL) images, framed as label-efficient learning of rare defects. The industry partner wanted a pipeline
that detects defects, grades severity and estimates power loss; the faculty advisor wanted a paper with a lightweight model (about 2M
parameters in the notes of the time, about 4M on the course poster). The data were BenchmarkEL \citep{pratt2023benchmark} (582
labelled cells, $512 \times 512$, 24 classes, and about 12k unlabelled cells on a public share) and UCF-EL \citep{fioresi2022ucfel}.

\paragraph{SA-UNet++ with Mean Teacher and SemiCE.} Phase~0 (the students' name for this period) ported SA-UNet++
(Gangopadhyay et al., IEEE Transactions on Emerging Topics in Computational Intelligence 9(2), 2025) to PyTorch and trained it with Mean Teacher
\citep{tarvainen2017meanteacher} and a consistency loss named SemiCE after PV-S3 \citep{jha2025pvs3}. The ``SA'' block computes
$\sigma(Q+K) \cdot V$, an elementwise sigmoid gate rather than dot-product attention; the model had 4,589,880 parameters, 88,000 of them
(1.9\%) in the gate blocks. The headline metric was defect mIoU, with gridline excluded as noisy. Three phase~0 results did not
survive. An ablation ``removing SA helps'' (0.2028 against 0.1792) and an SSL sweep (0.2580 at 20\% labels, 0.1323 at 10\%) predate a
fix to the noise augmentation: albumentations 2.x silently ignored the old \texttt{var\_limit} argument, giving noise of
$\sigma \approx 68$ instead of about 4 on the 0 to 255 scale, which erased cracks. An SSL gain of +7.4 points came from an 8$\times$
mismatch in optimiser steps, which produced the standing rule to match steps across arms. The square-root inverse-frequency class
weights (about 310:1) were later found to be computed over train, validation and test together. On 11 Oct the user recalled that
``SA SSL'' had raised IoU; the phase~1 notes traced this to a table that listed one arm, SA-SL at 20\% labels, seed 0, at two
checkpoints (0.338 at the final step, 0.301 at the best-validation step). In our reading it was a checkpoint artefact.

\paragraph{First reproductions and the poster.} Phase~0 also reproduced four papers (Fioresi, Pratt, PV-S3, SA-UNet++). Their code was
lost; the surviving records are the numbers used on the course poster (EE511 Computer Vision), ``Label-Efficient Segmentation of
Rare Defects in Electroluminescence Images of Solar Cells'' (Table~\ref{tab:repro-poster}). An earlier claim that SSL doubled UCF-EL
performance (10.83 to 21.19) compared two under-trained runs and is recorded as invalid. The poster's method figure had three panels: (a) a
U-Net of about 4M parameters with a pretrained light encoder and a full-resolution skip for 2~px cracks; (b) SSL inspired by ALHVR
\citep{lei2025alhvr} with M1 defect-aware sampling, M2 per-class thresholds $\gamma_k$ and M3 an EMA prototype memory; (c) severity
grading. Its presenting line was that at equal compute our SSL reimplementation showed no gain, that the published PV-S3 supervised
baseline (8.8) is under-trained, and that the PV-S3 code ignores background. By 6 Oct the phase~1 notes listed as outdated or wrong: the
Fig.~3 caption cited cross pseudo-supervision, which ALHVR does not use; M3 had been measured dead (Section~\ref{sec:repro-phase1}); the
step-matched control (0.491 $\pm$ 0.009) and the PV-S3 released-code finding were missing; and the SA-UNet++ row, the crack-width
wording and the move of all numbers to final-step values from the rebuilt pipeline were pending. The ``SSL 46.4'' record was later not reproduced
(Section~\ref{sec:repro-pvs3}).

\begin{table}[htbp]
\centering
\footnotesize
\setlength{\tabcolsep}{4pt}
\caption{Old era, lost pipeline: phase~0 reproduction records as used on the poster (percent, as recorded; code lost, numbers from
notes, never re-verified; test sets scored in every run; overlapping splits as later found). ``PV-S3 metric'': pooled IoU with every
pixel whose ground truth is background dropped.}
\label{tab:repro-poster}
\begin{tabular}{@{}p{0.31\linewidth}p{0.14\linewidth}p{0.25\linewidth}p{0.24\linewidth}@{}}
\toprule
Baseline & Published & Ours (as recorded) & Note \\
\midrule
Fioresi, supervised 100\%, 4-class / 5-class mean & 48.15\% / 57.30\% & 50.21\% / 59.42\% & 5-class includes background; 48.15 equals the IoU equivalent of the published F1 (our identification) \\
Pratt U-Net, per-image median, rule B & 13.33 to 16.00\% & 15.01\% & rule A 29.22\%, rule C 26.87\% \\
Jha, supervised 20\%, 30 epochs & 8.76\% & 13.13\% (PV-S3), 10.83\% (standard) & same under-trained regime \\
Jha, supervised 20\%, 12.0k steps & n/a & 56.61\% / 39.00\% & PV-S3 metric / standard \\
Jha, SSL 20\% (our reimplementation), 95.7k steps & 67.00\% & 58.85\% / 46.40\% & not reproduced in phase~1 \\
Supervised 20\%, step-matched, 95.8k steps & n/a & 60.79\% / 48.90\% & above our SSL \\
Jha, supervised 100\%, 59.8k steps & 70.84\% & 65.84\% / 50.21\% & PV-S3 metric / standard \\
\bottomrule
\end{tabular}
\end{table}

\paragraph{The severity-grading idea of the time.} Panel (c) and a ``Subpath A'' plan described post-processing of BenchmarkEL masks:
per-class connected components, a flag for components cut by the border, crack skeletons (length, medial-axis width, tortuosity), and
electrical isolation from a topology mask of classes $\{1,2,3,4,9\}$ in which crack pixels act as walls and components that touch no
busbar are flagged. Grades followed the crack modes of \citet{kontges2011risk}: A, a crack with no isolated area; B, an isolated area
that still glows; C, an isolated area that is dark. The notes held that power loss tracks the mode-C area rather than total crack
pixels, and that without paired current-voltage data a power estimate would rest on published correlations. Agreement between grades
from ground-truth and predicted masks was to be scored by Cohen's kappa. The instance stage existed with 57 unit tests; a compression
plan (pruning to about 2M parameters, INT8 with TensorRT on a Jetson) was never executed. Later severity work is in
Chapter~\ref{ch:deployment}.

\paragraph{The reset.} Before 2 Oct the user account on the single-GPU institute server (the second machine of later chapters) was
reset. The reproduction code, the UCF-EL validation split, the Subpath~A code, the label splits and the predictions were lost, and the
remote repository was two months stale, so all September work was unrecoverable.

\section{Phase 1: the rebuild (2 to 7 Oct 2026)}\label{sec:repro-phase1}

\paragraph{Setup check and foundation (2 Oct).} The server was found empty; the datasets and the ALHVR reference code were cloned and
11,993 cells of the external unlabelled pool downloaded. A check before any training found that the old loader expected colour masks
while the release ships index masks (24.6\% of pixels would have become ``ignore''); that the hard-coded class frequencies matched all
582 cells including test (on train only the weight ratio is 615:1, not 313:1); that a missing split file made the old code train on all
582 cells, test included; and that the noise augmentation fell back to a default about 10 to 15$\times$ stronger than intended. A new
repository was built with 96 tests: index masks read directly, missing splits an error, class weights from the labelled subset on train
only (weight 0 for absent classes), one metrics module with the \texttt{standard}, PV-S3 and per-image-median conventions, a plain
U-Net++ \citep{zhou2018unetpp} of 4,981,368 parameters trained from scratch, Adam at $10^{-3}$ with cosine decay to zero over 20,000
steps, and the rule ``final-step test is primary'' (best-validation selection on 54 cells flipped the sign of several comparisons).
Coupled augmentations from a shared random generator (both flips always together, blur always with noise) were fixed on 3 Oct, so the
pilots of 2 Oct are kept apart. An OS shutdown at 21:34 on 2 Oct killed the first ALHVR pilot at step 13,600, and reboots on 6 and 7
Oct killed queued runs; resume was built and tested bit-exact.

\paragraph{Arms and loss-off controls.} From 4 Oct a shared eight-GPU server (the training machine of later chapters) joined.
Mean Teacher used EMA 0.999, $\lambda$ ramped to 1.5 by step 4,000, our SemiCE and 4 labelled plus 4 unlabelled cells per step, with
the teacher evaluated. ALHVR \citep{lei2025alhvr} was ported following the reference code: two networks and optimisers, clean and
feature-perturbed passes, $\lambda_{\max} = 0.1$ ramped over the whole run. At 20\% labels (95 cells) the 383 other train cells were
the unlabelled pool. Table~\ref{tab:repro-phase1} gives the arms. Mean Teacher with $\lambda = 0$ (unlabelled images still forwarded)
scored 0.3653 $\pm$ 0.0183 against 0.3648 $\pm$ 0.0071 with $\lambda = 1.5$ (3 seeds each; paired $-0.000 \pm 0.024$, n.d.), and
teacher minus student was $+0.000 \pm 0.002$, since an EMA adds nothing when the learning rate ends at zero. ALHVR with its unlabelled
weight at 0 scored 0.3642 (paired on minus off $+0.005 \pm 0.023$, n.d.). The reading at the time was that the SSL arms' lead did not
come from their unlabelled losses; the leading hypothesis, batch-norm statistics over a half-unlabelled batch, was not shown.

\begin{table}[htbp]
\centering
\footnotesize
\caption{Old era, first pipeline (from-scratch U-Net++, original masks, weighted CE, batch 4, final step of 20,000 steps; the 50
test cells scored in every run; overlapping splits as later found; Mean Teacher with the invalid SemiCE). Test defect mIoU, standard,
mean $\pm$ sd over runs pooled from both machines (seeds 0 to 2 ran on both with the same initialisation, so $n$ overstates
independence). Welch 95\% intervals were computed after the runs, uncorrected; no rule was fixed before them.}
\label{tab:repro-phase1}
\begin{tabular}{@{}lrll@{}}
\toprule
Arm & $n$ & Test defect mIoU & Difference to SL-20 [95\% CI] \\
\midrule
SL-20 (supervised, 20\% labels) & 8 & 0.3435 $\pm$ 0.0329 & \\
SL-100 & 4 & 0.3383 $\pm$ 0.0060 & $-$0.0052 [$-$0.0330, +0.0226], n.d. \\
MT-20 (Mean Teacher, SemiCE) & 6 & 0.3662 $\pm$ 0.0061 & +0.0227 [$-$0.0050, +0.0504], n.d. \\
ALHVR-20 & 6 & 0.3724 $\pm$ 0.0103 & +0.0289 [+0.0007, +0.0570] \\
SA-SL-20 & 4 & 0.3399 $\pm$ 0.0104 & $-$0.0036 [$-$0.0324, +0.0251], n.d. \\
SA-SL-100 / SA-MT-20 & 1 / 1 & 0.3197 / 0.3384 & \\
Pratt U-Net, weight set A / C, 100\% & 1 / 1 & 0.3386 / 0.3149 & \\
\bottomrule
\end{tabular}
\end{table}

\paragraph{Label sweeps and diagnostics.} With labelled batch 8, supervised test defect mIoU was flat from 10\% to 100\% labels (0.357,
0.354 and 0.346, six runs each) and broke at 2\% (0.191, three runs); the two machines disagreed on SL-20 at batch 8 (means 0.346 and
0.362). Crack precision was 0.22 to 0.35 at recall 0.67 to 0.73 in every class-weighted run. A fit check found crack IoU 0.325 on the
47 training cells of a 10\% model (inactive 0.881 there, 0.417 on test): crack was not fitted even on training cells, the number
behind the user's worry of 7 Oct. A prototype diagnostic on four ALHVR runs found that the low-confidence region $\Omega_3$ held 71\%
of true crack pixels (119,384 of 168,511) on 383 held-out train cells, and that the other network's prototypes, which ALHVR's
prototype loss uses, labelled 4.6\% of them correctly; prototype modules were recorded as dead. On UCF-EL a supervised sweep
(DeepLabV3+ \citep{chen2018deeplabv3plus} with ResNet-50, 95,760 steps, 2,289 public test images, standard four-class mean) gave
0.5364 $\pm$ 0.0042 at 100\% (2 seeds), 0.4932 $\pm$ 0.0045 at 20\%, 0.4146 $\pm$ 0.0017 at 5\% and 0.2265 $\pm$ 0.0017 at 1\% (3
seeds each); the 1\% subset has no corrosion cell. Mean Teacher with SemiCE at 5\% and 1\% (7 to 8 Oct) scored below supervised
training (paired over 3 seeds, $-0.0913 \pm 0.0362$ and $-0.0742 \pm 0.0275$) and below its own $\lambda = 0$ controls; these runs test
SemiCE, not Mean Teacher.

\paragraph{Where phase 1 stood on 7 Oct.} The handoff of that evening summarised: on BenchmarkEL no unlabelled loss had shown an
effect, label count did not move the score between 5\% and 100\%, and crack could not be fitted on training cells; UCF-EL showed a clean
label gap, so the SSL question had moved there. While writing it, its author noticed that ALHVR's loss weight peaks only at the final
step, when the learning rate reaches zero (Section~\ref{sec:repro-alhvr}).

\section{Reproductions of published work}\label{sec:repro-papers}

Each subsection gives the claim and recipe, what we ran, our number in the paper's convention and in the standard one, and where the
gap came from; Table~\ref{tab:repro-summary} collects them. Runs before 7 Oct are old era (conditions above); P19 and the re-scorings
of 13 Oct used later code.

\subsection{Pratt et al.\ 2023: BenchmarkEL}\label{sec:repro-pratt}

The paper scores twelve models (four architectures, three class-weight sets) by per-class IoU per image, then the median over 50 test
cells chosen for a prevalence of cracks, gridline defects and inactive areas \citep{pratt2023benchmark}. Its headline, DeepLabv3+ with
custom weights (4c), has median crack 0.25, gridline 0.20, inactive 0.38 (mean 0.28). Weight set A is all 1; set C is background 0.2,
ribbons 3, gridline 10, inactive 10, crack 20, others 1. The paper states no optimiser, learning rate, epochs or batch size, does not
name the U-Net encoder, and does not say how a cell without the class enters the median.

On 3 Oct (old era) we trained a U-Net with an ImageNet VGG16 encoder (from the user's specification, not the paper; 23,750,424
parameters), sets A and C, flips only, and the pilot recipe (Adam $10^{-3}$, cosine, 20,000 steps, batch 4), one seed each on all 478
train cells, and scored three absent-class rules: A leaves the cell out, B scores 0, C scores 1 if nothing is predicted. Of the 50 test
cells, 50 contain crack, 37 gridline and 17 inactive. At the final step set A gave crack 0.149, gridline 0.339, inactive 0.420 under
rule A (mean 0.303; 0.145 under B, 0.449 under C); set C gave 0.176, 0.207, 0.445 (mean 0.276; 0.113 under B, 0.326 under C). In the
standard convention set A ranked first (mean of the three classes 0.331 against 0.284) while the paper ranks C first.

The report of the day took rule B as primary and called the paper ``approximately reproduced'' (set C 0.116 at its best-validation
step, set A 0.145, against the paper's U-Nets at 0.133 to 0.160). The audit analysis of 7 and 8 Oct showed that the paper cannot have used rule B:
with inactive in 17 of 50 test cells, rule B forces the inactive median to 0, yet 4c reports 0.38; and model 4a's inactive median 0.13 at
recall 0.15 rules out rule C. Under rule A, which we infer, no gap remains: set A 0.303 against the equal-weight U-Net 2a at 0.19,
the U-Net++ pilot arms (coupled-augmentation code, best-validation weights) at 0.280 to 0.295 against 4c at 0.28, and the U-Net++ arms'
crack medians 0.21 to 0.25 (the two Pratt U-Nets 0.149 and 0.176) against the paper's 0.13 to 0.26. The rule alone moves our mean about 3$\times$ (0.15 under B, 0.49 under C). On 13 Oct the same check on
supervised P-series models at 100\% labels (5 seeds, trained on revised masks, scored on the original masks, rule A) gave crack 0.354
$\pm$ 0.004, gridline 0.303 $\pm$ 0.010, inactive 0.443 $\pm$ 0.028 (mean 0.367 $\pm$ 0.012); in our reading not a like-for-like claim
(inferred rule, overlapping test cells, other training labels). Table~\ref{tab:T13b} has the final models.

\subsection{Fioresi et al.\ 2022: UCF-EL}\label{sec:repro-fioresi}

The paper trains DeepLabv3 with a COCO-pretrained ResNet-50, plain CE at learning rate 0.005, and calls a pixel background only if
$P(\text{background}) > 0.54$ \citep{fioresi2022ucfel}. It reports pixel-level F1 (headline, with 256 simulated images: crack 0.81,
contact 0.58, interconnect 0.33, corrosion 0.78) and an abstract mean IoU of 57.3\%, most likely a five-class mean with background (our
inference). One held-out set served as validation and test, and the public release (9,562 train, 2,289 test) is smaller than the
paper's 17,064 images. We trained DeepLabV3+ with an ImageNet ResNet-50 (shorter side 295, 256 crops, flips, SGD poly at
$5 \times 10^{-3}$, batch 16, final step, no checkpoint selection). On all labels for 100 epochs (59,760 steps) it scored 0.497 standard
(0.636 PV-S3), within the pre-set tolerance of 0.03 of the reference 0.4815, so the run with the paper's full augmentation was skipped
(0.481 is the mean of the IoU equivalents of the headline F1, IoU $=$ F1$/(2-$F1$)$ per class). Crack F1 of the 6 Oct supervised model
(100\%, 95,760 steps) was 0.806 against 0.81; contact, interconnect and corrosion reached 0.648, 0.498 and 0.795 against 0.58, 0.33 and
0.78. There is no deficit to explain; an early ``shortfall of 7.1'' came from setting a four-class mean beside the five-class 57.3\%.
It is not a strict reproduction: data, split, model and crops differ, and the paper's checkpoint was probably selected on its test set.

\subsection{Jha et al.\ 2025: PV-S3}\label{sec:repro-pvs3}

\paragraph{Claim and recipe.} Mean Teacher with ``Semi Cross-Entropy'' (SemiCE) handles class imbalance on EL images
\citep{jha2025pvs3}. On UCF-EL at 20\% labels the paper reports mean IoU 67.00\% for SSL against 8.76\% for supervised training (70.84\%
supervised on all labels). Recipe: DeepLabv3+ with ImageNet ResNet-50, 30 epochs, batch 16, SGD at $10^{-3}$, consistency weight 1.5,
SemiCE thresholds $(0.6, 0)$ in Tables 1 and 2 but $(0.2, 0.4)$ in its Section 4.2 and journal version. The metric is a pooled IoU over
the four defect classes; split and background treatment are not described.

\paragraph{The under-trained baseline.} In the PV-S3 convention (pooled, pixels whose ground truth is background dropped, so
background false positives vanish and recall is unchanged), our supervised model at 20\% labels scored 0.230 at the paper's budget (30
epochs at $10^{-3}$, 3,585 steps; 0.185 standard) and 0.559 at 12,000 steps at $5 \times 10^{-3}$ (0.448 standard), against 8.76\%;
the paper itself says its supervised model ``is not able to learn the features''. The step-matched control (95,760 steps, 3 seeds)
gave 0.491 $\pm$ 0.009 standard. The convention rewards false alarms: an all-crack predictor scores crack IoU 0.766 under it and 0.066
standard, and dilating predictions by 3~px raises the PV-S3 mean from 0.636 to 0.717 while the standard mean falls from 0.497 to 0.456.

\paragraph{Our reimplementation and the corrosion runaway.} Version 1 put the supervised loss on the strong view and collapsed
(Table~\ref{tab:repro-pvs3}); the strong view alone cost only 0.045 (0.448 to 0.403), so version 2 used the weak view and still
collapsed, predicting corrosion on 1,465,919 test pixels against 102,570 true. Measured without training, the final teacher labelled
0.96\% of unlabelled pixels as corrosion, 23$\times$ the labelled share and 16$\times$ the truth; 99.4\% of pixels cleared the 0.6
threshold, and the excess was still 21.6$\times$ at confidence 0.99. The PV-S3 convention hides it (corrosion IoU 0.067 standard, 0.879
PV-S3). Capping each class's pseudo-label share at 2$\times$ its labelled share repaired corrosion (IoU 0.631) but left the run 0.073
below the control; the smoothed cap was exceeded in 132 of 187 windows, one seed. Our version trained background, so it was not
faithful to the released code on that point, and the lost pipeline's ``SSL 46.4'' was not reproduced.

\begin{table}[htbp]
\centering
\footnotesize
\caption{Old era, first pipeline: PV-S3 reproduction on UCF-EL, 20\% labels (1,912 labelled images, our split lists; final step on the
2,289 public test images, scored on our evaluation grid with the shorter side at 295; UCF-EL has no validation split upstream, so
test was the only readout; 25 to 44\% of test cells have the same physical cell in train, as later found). Four-class mean IoU, pooled; ``PV-S3'' drops pixels
whose ground truth is background. DeepLabV3+ ResNet-50; the released code ran on the eight-GPU server, the rest on the single-GPU
server.}
\label{tab:repro-pvs3}
\begin{tabular}{@{}lrrr@{}}
\toprule
Run & Steps & Standard & PV-S3 \\
\midrule
Supervised, paper's budget (30 epochs, $10^{-3}$) & 3,585 & 0.185 & 0.230 \\
Supervised, $5 \times 10^{-3}$ & 12,000 & 0.448 & 0.559 \\
Supervised, step-matched (3 seeds) & 95,760 & 0.491 $\pm$ 0.009 & 0.607 $\pm$ 0.007 \\
Supervised, all 9,562 labels, 100 epochs & 59,760 & 0.497 & 0.636 \\
Reimplementation v1 (strong-view supervised loss) & 95,700 & 0.200 & 0.379 \\
Reimplementation v2, seed 0 / seed 1 & 95,700 & 0.276 / 0.283 & 0.568 / 0.436 \\
v2 with per-class pseudo-label cap (1 seed) & 95,700 & 0.418 & 0.580 \\
Released PV-S3 code as shipped (our splits) & 14,370 & 0.050 & 0.828 \\
\midrule
Published: supervised 20\% / SSL 20\% / supervised 100\% & & & 8.76\% / 67.00\% / 70.84\% \\
\bottomrule
\end{tabular}
\end{table}

\paragraph{The released code (5 Oct).} Run as shipped with our split lists (one student and two EMA teachers, 125.3M parameters;
learning rate $10^{-2}$ in its saved configuration against $10^{-3}$ in the paper; 14,370 steps), the code sets
\texttt{ignore\_index} to 0 in loss, loader, model and metric. UCF-EL encodes background as 0, so background never receives a gradient
and no test pixel was predicted as background. The standard mean was 0.050 (corrosion predicted on 3.33M pixels against 102,570, 32$\times$);
the PV-S3 convention gave 0.828 on our evaluation grid and 0.791 at native resolution, the code's own protocol. The code's own
end-of-training validation log (crack 0.9133, contact 0.7806, interconnect 0.5959, corrosion 0.8735) equals our native-resolution
per-class values to four decimals, so the value that matches the log is 0.791, not 0.828. Its supervised-only mode cannot run. Its
split lists name 16,327 images, 5,194 of them not in the public release, and put 1,769 of the 2,289 public test images in their
training lists and 437 in validation.

\paragraph{Where the gap came from.} The 8.76\% baseline is under-trained, and the convention adds to the gap. The SSL headline lives in
a convention blind to false positives on background, from code that never predicts background, with lists that put most public test
images in training; in the standard convention the step-matched control (0.491) is above every PV-S3 arm we ran. The ``SemiCE'' of our Mean Teacher
arms was not PV-S3's: the released loss is a softmax CE against the teacher's argmax where its maximum probability is at least 0.6; ours was a
per-class binary CE on raw logits with thresholds of half the class frequency (about 0.001 for crack); our PV-S3 reimplementation (v1, v2) used the softmax form. The paper's BenchmarkEL
comparison sets pooled IoU beside Pratt's per-image medians (Chapter~\ref{ch:related}). The journal version changed none of this.

\subsection{Gangopadhyay et al.\ 2025: SA-UNet++}\label{sec:repro-saunet}

On a UAV solar-panel defect set with four classes and no EL data, the paper reports Jaccard 0.95 $\pm$ 0.09 for SA-UNet++ against
0.92 $\pm$ 0.12 for U-Net++ (one run; spread over images; background probably in the mean, our inference)
(Gangopadhyay et al., 2025). Our corrected rerun on BenchmarkEL (old era, Table~\ref{tab:repro-phase1}) gave a
difference of $-0.0036$ [$-0.0324$, +0.0251] on test and $-0.0211$ [$-0.0510$, +0.0089] on validation plus test, both n.d. In our
reading the two do not conflict: the paper's gain is on large objects under a background-inclusive mean, and it never tested 2 to 3~px
structures.

\subsection{Lei et al.\ 2025: ALHVR (a port, not a reproduction of their numbers)}\label{sec:repro-alhvr}

ALHVR routes unlabelled pixels by the joint confidence of two networks and reports 81.59 rising to 90.56 Dice on cardiac MRI with
10\% labels \citep{lei2025alhvr}; it does not target rare classes. Our port differed in optimiser (Adam with cosine to zero against
SGD at a constant rate), supervised loss (weighted CE against CE plus Dice), batch (4 plus 4 against 8 plus 8), classes (24 against 4)
and checkpoint (final step against best validation). The audit of 7 and 8 Oct found that the learning-rate-weighted mean of $\lambda$ was
0.0156 against 0.0393 in the reference, that only 8.1\% of our learning-rate mass fell where $\lambda \geq 0.05$ (37\% of the
reference's iterations), and that the logged ratio of weighted unlabelled to supervised loss rose from 0.7\% at step 2,000 to 89.4\%
at step 19,000, when the learning rate was 0.6\% of its peak. The loss-off result ($+0.005 \pm 0.023$) was close to null by
construction and showed neither a gain nor its absence; no wiring error was found. The fair port of 23 Oct is row AL1 of
Table~\ref{tab:T21}.

\subsection{Hanifi et al.\ 2025: P19 (14 Oct)}\label{sec:repro-p19}

\citet{hanifi2025multiclass} train six encoder-decoder models for 29-class segmentation on BenchmarkEL \texttt{20221008} and report
per-class IoU pooled over test cells, with group means over feature, defect and all classes. At the user's request, protocol P19 was
locked on 14 Oct before any run, and the authors' released code ran unchanged at a fixed commit, with their TensorFlow and Keras
versions, on the eight-GPU server that day: release split 553 / 70 / 72, no augmentation, input $416 \times 608$, weighted CE with
their weights, Adam $10^{-3}$, batch 8, early stopping on validation loss. Six models ran at seed 0, and the headline VGG16-UNet also
at seeds 1 and 2. Trainable parameters equal the paper's for all six models, and our per-class counts reproduce their evaluation
function. Table~\ref{tab:repro-p19} uses their convention.

\begin{table}[htbp]
\centering
\footnotesize
\caption{P19 (14 Oct; locked protocol; not old era): VGG16-UNet of \citet{hanifi2025multiclass}, ours mean $\pm$ sd over 3 seeds
against the paper's single run. Their convention: per-class IoU pooled over the 72 release test cells at $208 \times 304$, an absent
class scoring 0 in group means. Not comparable with any number of this thesis (other split, labels, resolution and recipe).}
\label{tab:repro-p19}
\begin{adjustbox}{max width=\linewidth}
\begin{tabular}{@{}lll@{}}
\toprule
Quantity & Ours (3 seeds) & Paper (1 run) \\
\midrule
Mean IoU, 12 feature classes & 0.678 $\pm$ 0.020 & 0.663 \\
Mean IoU, 16 defect classes & 0.246 $\pm$ 0.018 & 0.326 \\
Mean IoU, the 12 defect classes present in our test & 0.328 $\pm$ 0.024 & 0.369 \\
Mean IoU, 29 classes & 0.416 $\pm$ 0.015 & 0.494 \\
Frequency-weighted IoU & 0.884 $\pm$ 0.011 & 0.879 \\
Crack / inactive / gridline & 0.352 $\pm$ 0.008 / 0.469 $\pm$ 0.032 / 0.316 $\pm$ 0.028 & 0.335 / 0.550 / 0.297 \\
\bottomrule
\end{tabular}
\end{adjustbox}
\end{table}

The feature means, crack, gridline and the frequency-weighted IoU reproduce. The 16-class defect mean is lower mostly because our test
set has no dead-cell or belt-mark cell, which score 0 here against the paper's 0.502 and 0.293; over the 12 classes present in both,
half the remaining gap is the rings class (0.017 against 0.285). Inactive is lower (our test has 9 inactive cells). The paper's model
ranking does not reproduce (MobileNet-UNet, 0.353 on present defect classes, is level with VGG16-UNet; the paper's margin of 0.031 is
under two seed standard deviations of 0.018), and its 29-class mean cannot be recomputed from its own per-class figure (0.471 over 28
classes, 0.455 over 29, against 0.494). Its 551 / 67 / 77 re-split is not released and its 4$\times$ augmentation is not described;
neither like-for-like run (their model on our split, or ours on theirs) was made.

\subsection{Re-scoring in the SEiPV-Net convention (13 Oct)}\label{sec:repro-seipv}

SEiPV-Net \citep{eesaar2023seipv} reports mIoU 0.8604, F1 0.9426, Dice 0.8531 and IoU 0.7447 on the 50 legacy test cells, over 24
classes, with no per-class numbers, no definition of the averaging and no code. On 13 Oct we re-scored 18 final legacy-split P-series
checkpoints (Chapter~\ref{ch:legacy}) on the original \texttt{20211104} masks under several rules (read-only, one confusion matrix per
test cell). For the supervised 100\% models (5 seeds) the class mean over classes present in test ground truth was 0.514, the
pixel-weighted F1 0.933, and the per-image mean IoU with absent-and-unpredicted classes scored 1 was 0.840. In our reading (descriptive;
their rule is not confirmed; the 50 cells overlap our training pool) SEiPV-Net's numbers resemble such background-dominated quantities: our pixel-weighted F1 moves only 0.012
between 5\% and 100\% labels. Table~\ref{tab:T13b} repeats this for the final models.

\begin{table}[htbp]
\centering
\footnotesize
\caption{Reproduction summary. Rows from before 7 Oct are old era (first pipeline; test sets scored in every run; overlapping splits as
later found); P19 and SEiPV used later code. Published percentages as printed.}
\label{tab:repro-summary}
\begin{adjustbox}{max width=\linewidth}
\begin{tabular}{@{}p{0.21\linewidth}p{0.14\linewidth}p{0.29\linewidth}p{0.30\linewidth}@{}}
\toprule
Paper and number & Published & Ours (convention) & Explained by \\
\midrule
Pratt 4c, mean of 3 medians & 0.28 & set A 0.303; U-Net++ pilots 0.280 to 0.295 (rule A) & absent-class rule (A inferred, not B); no gap \\
Fioresi, crack F1 & 0.81 & 0.806 (supervised 100\%) & no deficit; data, split, model, selection differ \\
PV-S3 supervised 20\% & 8.76\% & 0.230 at their budget, 0.559 at 12k steps (PV-S3) & under-trained baseline \\
PV-S3 SSL 20\% & 67.00\% & v2 0.568 / 0.276; released code 0.828 / 0.050 (PV-S3 / standard) & convention blind to background false positives; background never predicted; test images in training lists \\
SA-UNet++ gate & 0.95 vs 0.92 Jaccard & $-0.0036$, n.d. (standard) & other domain and object scale \\
ALHVR, Dice at 10\% & 81.59 to 90.56 & loss-off $+0.005 \pm 0.023$, n.d. & loss barely applied; not a fair test \\
Hanifi VGG16-UNet, crack & 0.335 & 0.352 $\pm$ 0.008 (theirs, 3 seeds) & reproduces; defect means lower from absent test classes \\
SEiPV-Net mIoU / F1 & 0.8604 / 0.9426 & 0.840 / 0.933 (per-image, absent $=1$ / pixel-weighted F1) & undefined averaging; rule not confirmed \\
\bottomrule
\end{tabular}
\end{adjustbox}
\end{table}

\section{The audit and reset of 7 and 8 October}\label{sec:repro-audit}

At 23:37 on 7 Oct the user asked for the handoff to be read and wrote (typing errors corrected, no quotation) that the model had been
trained on 47 images and tested on the same 47 to reach 0.3 IoU on cracks, that something was definitely wrong, that the goal remained
a method inspired by ALHVR that focuses on the rare classes, and that they suspected minor training misconfigurations behind the zero
effect of the unlabelled losses. They asked for
analysis of the datasets, models, issues and published papers before any new run. The audit of 7 and 8 Oct found:
\begin{enumerate}
  \item \textbf{The test set was scored in every run and used for choices.} The final-step test number was the primary readout, and
    the comparisons of Table~\ref{tab:repro-phase1} were Welch tests on the 50 test cells (38 tests, uncorrected; one survives
    Bonferroni). A claim that one supervised budget was best had already been withdrawn as selection on test. UCF-EL has no
    validation split upstream, and Fioresi's official code validates on the test set.
  \item \textbf{Overlaps.} 20 of 54 validation and 24 of 50 test cells came from modules with train cells. The external pool contained
    55 labelled cells (48 train, 6 validation, 1 test), removed by exact match only, and kept 196 cells from validation modules and 88
    from test modules. The first class weights used test labels. On UCF-EL the carved validation split sat inside the unlabelled pool,
    and 25 to 44\% of test cells have the same physical cell, at another stress stage, in train. Section~\ref{sec:data-split} has the
    later module audit.
  \item \textbf{Metric issues.} Rule B was the primary Pratt rule though the paper's numbers exclude it; the PV-S3 convention hid a
    runaway class; under class weights defect mIoU largely measured crack over-prediction. Two independent masks of one duplicated
    training image agree at crack IoU 0.257, and a perfect mask shifted by 1~px scores 0.47 to 0.56, so strict crack-IoU differences
    of $\pm 0.03$ could not be interpreted (Chapter~\ref{ch:data} measures this on the later label sets).
  \item \textbf{Our SemiCE was neither PV-S3's nor a valid loss.} Protocol P2, locked before it ran at 00:44 on 8 Oct, probed the final
    Mean Teacher checkpoint: 100\% of unlabelled pixels had two or more ``positive'' classes (background on 98.9\%, crack on 20.3\%
    against a true share of about 0.19\%). The gradient was real (0.25 $\pm$ 0.35 of the supervised gradient) but nearly orthogonal to
    it (cosine 0.07). Every Mean Teacher result of the period, on both datasets, tested this loss, not Mean Teacher. P2 also found that a
    correct confidence rule (maximum probability $\geq 0.95$) would give crack 0.024\% of the confident pixels, about 8$\times$ under
    its labelled share: the rare-class problem of pseudo-labels, measured on our data.
  \item \textbf{ALHVR was barely applied} (Section~\ref{sec:repro-alhvr}).
  \item \textbf{The lost implementation.} Phase~0 results could not be re-run, ``SSL 46.4'' was not reproduced, and the old label
    subsets could not be recovered because their split files had never been committed.
  \item \textbf{The recipe, not a bug.} Protocol P1 (overfit check, locked before running, 8 Oct) trained on 10 crack-containing train
    cells for 3,000 steps and scored the same cells: crack IoU 0.612 with weighted CE and no augmentation, 0.861 with CE plus soft Dice,
    0.381 with the training augmentation; masks were aligned exactly, and bf16 matched fp32 over 3 seeds (0.668 $\pm$ 0.031 against
    0.656 $\pm$ 0.048). The training recipe, mainly blur and $\pm 5^\circ$ rotation on crack labels a median 1.75~px wide and
    secondly the loss, not a code or data bug, explained ``0.3 on its own training images''.
  \item \textbf{Design.} One labelled subset for every run, seeds pooled across machines with shared initialisation, and class weights
    that changed with the subset, so label fractions were compared at different operating points.
\end{enumerate}

\paragraph{What the analysis-first restart changed.} The research repository was started at 23:59 on 7 Oct with the P1 and P2
protocols committed before either ran; from then on every experiment had a protocol committed before its runs, and anything outside its
rules was labelled exploratory. On 8 Oct a literature survey of 69 verified sources, first-hand analyses of both datasets and a
situation report followed. The report concluded that neither SSL method had actually been tested, that strict crack IoU on BenchmarkEL
already sat at the level of label noise, and that the field's evaluation was weak (hand-picked test sets or test sets reused for validation,
conventions that drop background or take per-image medians, test images in training lists, under-trained baselines, single runs). The armed external-pool Mean Teacher
queue, which would have tested SemiCE, was held at 01:18 on 8 Oct and never ran. The rules kept were: match optimiser steps, report
final-step weights, never compare arms across machines, at least three seeds, one change per arm, and old-era numbers only in their own
tables. UCF-EL was set aside for label-efficiency claims until a module-disjoint re-split, which was never made. The P-series that began
on 10 Oct, on a corrected supervised recipe and revised labels, is Chapter~\ref{ch:legacy}; it still scored the 50 test cells in every
run (Table~\ref{tab:T21}), which the benchmark protocol replaced (Chapter~\ref{ch:protocol}).
